% Pursuing Stacks (cotes 134-2 to 134-8): archive cote 134-5, site page 67 (PDF page 68), typescript page 254.
% typescript. Draft compared: ps4.tex:3377-3441.
% AI-transcribed by Claude Opus 5.5 (claude-opus-5-5, adaptive thinking, effort high), Message Batches API, batch ps2 (msgbatch_01NVwr1Yuo8qJf5rykeGBXPK), 2026-11-12; not yet checked by eye.
\scanpage{134-5/67}
\mspage{254}

This array of immediate corollaries of th.1 do convince me that this
statement is indeed ``the'' natural generalization of the ``awkward'' main
result of section 51. All we have to do is to prove theorem 1 then.

The three conditions of theorem 1 can be rewritten simply as
\[ \text{(i)} \qquad i^{*}(C) \subset A^{\wedge}_{c} \qquad \text{(using H 4')} \ , \]
\[ \text{(ii)} \qquad i^{*}(M_{as}) \subset A^{\wedge}_{loc.as} \ , \]
\[ \text{(iii)} \qquad i^{*}(C) \subset A^{\wedge}_{loc.as} \ , \]
and because of
\[ C \subset M_{as} \ , \quad A^{\wedge}_{c} \subset A^{\wedge}_{loc.\,as} \ , \]
it follows tautologically that (i) and (ii) both imply (iii). On the
other hand, (iii)$\Rightarrow$(i) by criterion H 3'), and the definition of the
homotopy \inserted{interval} structure of $A^{\wedge}$ in terms of \xout{homotopy interval} $A^{\wedge}_{loc.as}$ . Thus
we are left with proving (i)$\Rightarrow$(ii). But (i) can be rewritten as
\[ i^{*}(M_{c}) \subset A^{\wedge}_{c} \ , \]
which implies
\[ (*) \qquad i^{*}(M_{c}) \subset A^{\wedge}_{loc.as} \ . \]
\deleted{To prove} \over{t}{T}hat the latter condition implies (ii)\deleted{, replacing A by any lo-%
calization, we are reduced to the case when A has a final object,
hence is aspheric, and to prove in this case that (*) implies}
\deleted{\[ (**) \qquad i^{*}(M_{as}) \quad A^{\wedge}_{as} \ , \]}
\deleted{i.e. that i is $M_{as}$-aspheric. On the other hand, because of our assump-%
tion on A, $A^{\wedge}_{loc.as}$ \quad $A^{\wedge}_{as}$ , hence (*) implies}
\deleted{\[ (*') \qquad i^{*}(M_{c}) \quad A^{\wedge}_{as} \ . \]}
\deleted{The conclusion} now follows from the \inserted{(corollary to the)} following tautology \inserted{(with $N=M_c$)}, which should
have been stated as a corollary to prop. 1 above (p. 247):

\uline{Lemma}: Let $(M,M_{as})$ be an asphericity structure, generated by the full
subcategory N of M, let A be a small category, and i: A$\rightarrow$M a functor
\uline{factoring through} N . Then \xout{for} i is $M_{as}$-aspheric iff for any x in N ,
$i^{*}(x)$ is an aspheric object of $A^{\wedge}$ .

\uline{Corollary} The functor i is locally \inserted{($M_{as}$-)} aspheric (resp. totally $M_{as}$-aspheric)
iff \xout{for} $i^{*}(N) \subset A^{\wedge}_{loc.as}$ (resp. $i^{*}(N) \subset A^{\wedge}_{tot.as}$ ).

\margin{67}
